% Appendix: what the five datasets actually are, and how each was constructed.
% Written out because two of the five are subsets or reformulations rather than
% the benchmark a reader would assume from the name.
\section{Datasets and methodology}
\label{app:datasets}

\subsection{The five datasets}

Every dataset is reduced to a set of standalone declarative claims with a binary
truth label, so that a single contrast pair can be built for each. Where the
source is a question-answering or entailment task, the reformulation is stated
below. Table~\ref{tab:datasets} summarizes the result.

\begin{table}[h]
\centering
\small
\begin{tabular}{llrrl}
\toprule
Dataset & Source & Items & Median tokens & Claim form \\
\midrule
Geometry of Truth & \citet{marks2024geometry} & $5030$ & $11$ & curated statement \\
TruthfulQA        & \citet{lin2022truthfulqa} & $1634$ & $22$ & question $+$ candidate answer \\
RTE               & \citet{wang2018glue}      & $277$  & $58$ & premise $+$ hypothesis \\
BoolQ             & \citet{clark2019boolq}    & $3270$ & $121$ & passage $+$ yes/no question \\
IMDB              & \citet{maas2011imdb}      & $5000$ & $187$ & review $+$ sentiment claim \\
\bottomrule
\end{tabular}
\caption{The five datasets after reduction to labelled claims. Token counts are
GPT-2 tokens of the claim text, before the contrast-pair suffix is appended.
IMDB is capped at $5000$ items, balanced across labels; the others are used in
full.}
\label{tab:datasets}
\end{table}

\paragraph{Geometry of Truth.}
We use four of the fifteen topic files released with
\citet{marks2024geometry}: \texttt{cities}, \texttt{sp\_en\_trans},
\texttt{larger\_than} and \texttt{companies\_true\_false}, read from the
authors' repository. These are short affirmative statements of a single fact,
for example \emph{``The city of Krasnodar is in Russia.''} or \emph{``Fifty-one
is larger than fifty-two.''}

The eleven files we do not use include the negated counterparts of three of the
four we do (\texttt{neg\_cities}, \texttt{neg\_sp\_en\_trans},
\texttt{smaller\_than}), the logical compositions
(\texttt{cities\_cities\_conj}, \texttt{cities\_cities\_disj}), and the
uncurated sets (\texttt{common\_claim\_true\_false},
\texttt{counterfact\_true\_false}). Our subset is therefore the affirmative,
templated portion of the benchmark, which is reflected in its statistics: it has
the shortest claims of the five datasets by a wide margin and the least
variation in length.

This matters for interpreting the results, because Geometry of Truth is the
dataset on which detection is strongest and transfers best
(Section~\ref{sec:indomain}). Whether that reflects the factual-recall content
of the claims or their surface uniformity is not settled by the present
experiments. Adding the negated files would separate the two, since they hold
the same facts in a harder syntactic form, and we flag this as the most
informative single extension to the data.

\paragraph{TruthfulQA.}
Each question is paired with one candidate answer to form a claim, labelled by
whether that answer is among the correct ones. Questions are kept together under
a shared group id so that a question's answers cannot be split across the train
and test halves.

\paragraph{RTE.}
The premise and hypothesis are concatenated into a single claim, labelled by
whether the premise entails the hypothesis. It is the smallest dataset at $277$
items, which is why several of its cells fall below the thresholds used in
Sections~\ref{sec:indomain} and the detection analysis.

\paragraph{BoolQ.}
The passage and its yes/no question are concatenated, labelled by the answer.
Claims are long and highly variable in length.

\paragraph{IMDB.}
A review is paired with a claim about its sentiment, labelled by whether the
claim matches the review's label. Capped at $5000$ items with equal numbers of
each label; the other four datasets are used in full.

\subsection{Splitting and group integrity}

Every dataset is split once into train and test halves by a fixed seed, and the
same split is reused by every model, so that a pair of models is always compared
on identical items. Items sharing a group id --- the answers to one TruthfulQA
question, for instance --- are assigned to the same half, so that a probe cannot
be fitted on one answer to a question and scored on another. Truncation to a
cap is stratified by label and respects groups.

\subsection{Thresholds}

False agreement is a rare event on some datasets, and a rate estimated from a
handful of items is not informative. Two thresholds are therefore applied and
stated wherever they affect a number. Cells with fewer than $20$ false
agreements in the evaluated split are excluded from detection results; this
removes $165$ of $611$ cells in the in-domain analysis, falling hardest on RTE,
which retains $18$. No threshold is applied to the false-agreement rates
themselves, which are averaged over all cells.
